% ═══════════════════════════════════════════════════════════════
% LEHRERHINWEISE: Der Gleichstrommotor (Kl. 9, DS 09a)
% Basiert auf: templates/mint/lehrerhinweise-mint.tex
% ═══════════════════════════════════════════════════════════════

\documentclass[11pt, a4paper]{article}

\usepackage[utf8]{inputenc}
\usepackage[T1]{fontenc}
\usepackage[ngerman]{babel}

\usepackage{helvet}
\renewcommand{\familydefault}{\sfdefault}

\usepackage[margin=2cm, top=2cm, bottom=2cm]{geometry}
\usepackage{amsmath, amssymb}
\usepackage{siunitx}
\sisetup{locale = DE, per-mode = symbol}

\usepackage{enumitem}
\usepackage{tabularx}
\usepackage{booktabs}
\usepackage{array}
\usepackage{longtable}

\usepackage{xcolor}
\usepackage{tcolorbox}
\tcbuselibrary{skins, breakable}

\usepackage{fancyhdr}
\usepackage{lastpage}
\usepackage{hyperref}

\setlist{nosep, leftmargin=1.5em}

% ═══════════════════════════════════════════════════════════════
% FARBEN
% ═══════════════════════════════════════════════════════════════

\definecolor{physik}{HTML}{2c5aa0}
\definecolor{hellgrau}{HTML}{F0F0F0}
\definecolor{grau}{HTML}{666666}
\definecolor{success}{HTML}{2a7a4b}
\definecolor{warning}{HTML}{b35c00}
\definecolor{error}{HTML}{c62828}

\definecolor{basis}{HTML}{2a7a4b}
\definecolor{standard}{HTML}{2c5aa0}
\definecolor{erweiterung}{HTML}{7b1fa2}

% ═══════════════════════════════════════════════════════════════
% VARIABLEN
% ═══════════════════════════════════════════════════════════════

\newcommand{\fachid}{physik}
\newcommand{\fach}{Physik}
\newcommand{\thema}{Der Gleichstrommotor}
\newcommand{\klasse}{9}
\newcommand{\stession}{09a}
\newcommand{\stundenanzahl}{--}
\newcommand{\reihe}{Motor und Induktion}
\newcommand{\dauer}{40 Min.}
\newcommand{\lerngruppengroesse}{24}

\colorlet{fachfarbe}{physik}

% ═══════════════════════════════════════════════════════════════
% KOPF-/FUSSZEILE
% ═══════════════════════════════════════════════════════════════

\pagestyle{fancy}
\fancyhf{}
\renewcommand{\headrulewidth}{0.4pt}
\lhead{\small\textcolor{fachfarbe}{\textbf{\fach{} Kl. \klasse}} -- \thema{} -- \textbf{Lehrerhinweise}}
\rhead{\small LH-\stession}
\cfoot{\small\thepage/\pageref{LastPage}}

% ═══════════════════════════════════════════════════════════════
% BEFEHLE
% ═══════════════════════════════════════════════════════════════

\newcommand{\Gkurs}{\textcolor{basis}{\textbf{[G]}}}
\newcommand{\Ekurs}{\textcolor{erweiterung}{\textbf{[E]}}}
\newcommand{\lsg}[1]{\textcolor{error}{\textbf{#1}}}

% ═══════════════════════════════════════════════════════════════
% BOXEN
% ═══════════════════════════════════════════════════════════════

\newtcolorbox{infobox}[1][]{
    colback=hellgrau,
    colframe=fachfarbe,
    boxrule=1pt,
    arc=0pt,
    left=8pt, right=8pt, top=6pt, bottom=6pt,
    fonttitle=\bfseries,
    title=#1,
    breakable
}

\newtcolorbox{warnbox}[1][]{
    colback=white,
    colframe=warning,
    boxrule=1pt,
    arc=0pt,
    left=8pt, right=8pt, top=6pt, bottom=6pt,
    fonttitle=\bfseries,
    title=#1,
    breakable
}

\newtcolorbox{successbox}[1][]{
    colback=white,
    colframe=success,
    boxrule=1pt,
    arc=0pt,
    left=8pt, right=8pt, top=6pt, bottom=6pt,
    fonttitle=\bfseries,
    title=#1,
    breakable
}

% ═══════════════════════════════════════════════════════════════
% DOKUMENT
% ═══════════════════════════════════════════════════════════════

\begin{document}

% ═══════════════════════════════════════════════════════════════
% TITEL
% ═══════════════════════════════════════════════════════════════

\begin{center}
{\Large\textbf{\textcolor{fachfarbe}{Lehrerhinweise: \thema}}}\\[0.3em]
{\normalsize \fach{} Klasse \klasse{} -- Stunde \stession{} -- Reihe ``\reihe''}
\end{center}

\vspace{1em}

% ═══════════════════════════════════════════════════════════════
% 1. KURZÜBERSICHT
% ═══════════════════════════════════════════════════════════════

\section*{1. Kurzübersicht}

\begin{infobox}
\begin{tabular}{ll}
\textbf{Thema:} & \thema{} \\
\textbf{Klasse:} & \klasse{} (G-Kurs / E-Kurs) \\
\textbf{Zeit:} & \dauer \\
\textbf{Lernziele:} & Aufbau, Funktionsprinzip, Drehrichtung ändern \\
\textbf{Methoden:} & Demonstration, Lernpfad, Minitest \\
\textbf{Material:} & AB-\stession, LP-\stession, MT-\stession, Motormodell \\
\end{tabular}
\end{infobox}

% ═══════════════════════════════════════════════════════════════
% KURZVERSION
% ═══════════════════════════════════════════════════════════════

\begin{warnbox}[Kurzversion 25 Min (z.B. bei Schülervortrag)]
\textbf{Nur LP + LB, kein MT:}
\begin{itemize}
    \item 3' Einstieg: Motormodell zeigen
    \item 15' LP-\stession{} durcharbeiten (Schritte 1--4, Simulation!)
    \item 5' AB-\stession: Merkkasten 1 + Aufgabe 1 (Lückentext)
    \item 2' Sicherung: Was ist der Kommutator?
\end{itemize}
\textbf{AB-\stession{} mitnehmen} -- Rest als Hausaufgabe oder nächste Stunde.
\end{warnbox}

\vspace{1em}

% ═══════════════════════════════════════════════════════════════
% 2. STUNDENVERLAUF
% ═══════════════════════════════════════════════════════════════

\section*{2. Stundenverlauf (Vollversion mit MT)}

\begin{longtable}{|c|p{2cm}|p{7cm}|p{3.5cm}|}
\hline
\textbf{Min.} & \textbf{Phase} & \textbf{Inhalt} & \textbf{Material} \\
\hline
\endhead

0--3 & Einstieg & DE: Motormodell zeigen, laufen lassen & Modell (groß/klein) \\
\hline

3--15 & Erarbeitung & LP-\stession{} durcharbeiten (Simulation!) & Tablets, QR-Code \\
\hline

15--20 & Sicherung & Kurze Besprechung: Was ist der Kommutator? & Plenum \\
\hline

20--38 & \textbf{Leistung} & \textbf{MT-\stession: Gleichstrommotor (25 BE)} & Tablets, QR-Code \\
\hline

38--40 & Abschluss & Ergebnis notieren, Ausblick Induktion & --- \\
\hline
\end{longtable}

% ═══════════════════════════════════════════════════════════════
% 3. DEMONSTRATIONSEXPERIMENT
% ═══════════════════════════════════════════════════════════════

\section*{3. Demonstrationsexperiment (Pflicht lt. Rahmenplan)}

\begin{infobox}[DE: Gleichstrommotor]
\textbf{Variante A: Motormodell}
\begin{itemize}
    \item Modell mit sichtbarem Kommutator verwenden
    \item Polung umkehren $\rightarrow$ Drehrichtung ändert sich
    \item Magnete umdrehen $\rightarrow$ Drehrichtung ändert sich
    \item Beides ändern $\rightarrow$ Drehrichtung bleibt gleich!
\end{itemize}

\textbf{Variante B: Selbstbau-Motor (wenn Zeit)}
\begin{itemize}
    \item Kupferdraht-Spule auf Achse
    \item Neodym-Magnete, Batterie
    \item Büroklammern als Lager
\end{itemize}
\end{infobox}

% ═══════════════════════════════════════════════════════════════
% 4. EXTERNE RESSOURCEN
% ═══════════════════════════════════════════════════════════════

\section*{4. Alternative Simulationen / Externe Ressourcen}

\begin{warnbox}[Empfohlene externe Ressourcen]
\textbf{LEIFI Physik} (sehr empfohlen!)
\begin{itemize}
    \item \url{https://www.leifiphysik.de/elektrizitaetslehre/elektromotoren-und-generatoren}
    \item Animation zum Gleichstrommotor mit Kommutator
    \item Gute Grafiken zum Ausdrucken/Projizieren
\end{itemize}

\textbf{PhET Colorado}
\begin{itemize}
    \item \url{https://phet.colorado.edu/de/simulations/faradays-law}
    \item ``Faraday's Law'' -- zeigt Induktion (gut für nächste Stunde)
\end{itemize}

\textbf{YouTube (zur Vorbereitung)}
\begin{itemize}
    \item SimpleClub: ``Gleichstrommotor einfach erklärt''
    \item musstewissen Physik: ``Wie funktioniert ein Elektromotor?''
\end{itemize}

\textbf{Eigene Simulation:} SIM-\stession-gleichstrommotor.html
\end{warnbox}

% ═══════════════════════════════════════════════════════════════
% 5. DIFFERENZIERUNG
% ═══════════════════════════════════════════════════════════════

\section*{5. Differenzierung}

\begin{tabular}{|l|p{4cm}|p{4cm}|p{4cm}|}
\hline
& \Gkurs{} \textbf{G-Kurs (Basis)} & \textbf{Standard} & \Ekurs{} \textbf{E-Kurs (Erw.)} \\
\hline
\textbf{Ziel} & Bauteile benennen, Funktion beschreiben & Lorentzkraft erklären & Transfer, Wechselstrom \\
\hline
\textbf{MT-Aufgaben} & 1, 2, 4 (13 BE) & Alle (25 BE) & Alle + Zusatzfrage \\
\hline
\textbf{Hilfen} & Begriffskarten, Simulation langsam & LP vollständig & Nur Impulse \\
\hline
\end{tabular}

\vspace{1em}

\textbf{Zusatzfrage E-Kurs:} ``Warum braucht man bei Wechselstrom keinen Kommutator?''

% ═══════════════════════════════════════════════════════════════
% 6. HÄUFIGE SCHÜLERFEHLER
% ═══════════════════════════════════════════════════════════════

\section*{6. Häufige Schülerfehler}

\begin{warnbox}[Typische Fehler]
\begin{tabular}{|p{4.5cm}|p{9cm}|}
\hline
\textbf{Fehler} & \textbf{Reaktion / Hilfe} \\
\hline
Verwechslung Rotor/Stator & Eselsbrücke: ``Rotor \textbf{r}otiert, Stator \textbf{st}eht still'' \\
\hline
Kommutator-Funktion unklar & Animation in Simulation zeigen, Zeitlupe nutzen \\
\hline
Rechte-Hand-Regel falsch & Daumen = I (techn.), Zeigefinger = B, Mittelfinger = F \\
\hline
Drehrichtung ändern & ``Nur EINE Sache ändern, nicht beide!'' \\
\hline
\end{tabular}
\end{warnbox}

% ═══════════════════════════════════════════════════════════════
% 7. MATERIAL-CHECKLISTE
% ═══════════════════════════════════════════════════════════════

\section*{7. Material-Checkliste}

\begin{itemize}
    \item[$\Box$] Motormodell (groß oder klein) mit sichtbarem Kommutator
    \item[$\Box$] Netzgerät oder Batterie
    \item[$\Box$] QR-Code: LP-\stession{} (Projektion oder Ausdruck)
    \item[$\Box$] QR-Code: MT-\stession{} (Projektion)
    \item[$\Box$] Tablets/Smartphones für SuS
    \item[$\Box$] AB-\stession.pdf kopiert (\lerngruppengroesse{} Kopien)
    \item[$\Box$] Optional: LEIFI-Animation am Beamer
\end{itemize}

% ═══════════════════════════════════════════════════════════════
% 8. LÖSUNGEN MT-09a
% ═══════════════════════════════════════════════════════════════

\section*{8. Lösungen MT-\stession{} (25 BE)}

\begin{successbox}[Musterlösungen]

\textbf{Aufgabe 1: Zuordnung (5 BE)}
\begin{itemize}
    \item (A) = Rotor (Spule)
    \item (B) = Stator (Magnet)
    \item (C) = Bürsten
    \item (D) = Kommutator
    \item (E) = Drehachse
\end{itemize}

\textbf{Aufgabe 2: Warum dreht sich die Spule? (4 BE)}\\
\lsg{(b)} Die Lorentzkraft wirkt auf die stromdurchflossenen Leiter der Spule.

\textbf{Aufgabe 3: Kraftrichtung (4 BE)}\\
\lsg{(a)} Nach oben ($\uparrow$).\\
\textit{Begründung: Rechte-Hand-Regel: Daumen zu dir (I $\odot$), Zeigefinger $\rightarrow$ (B), Mittelfinger zeigt nach oben (F $\uparrow$).}

\textbf{Aufgabe 4: Kommutator-Aufgabe (4 BE)}\\
\lsg{(b)} Er kehrt die Stromrichtung in der Spule bei jeder halben Umdrehung um.

\textbf{Aufgabe 5: Drehrichtung ändern (4 BE)}\\
\lsg{(b)} Durch Umpolen der Spannung ODER Umpolen des Magnetfelds.

\textbf{Aufgabe 6: Transfer (4 BE)}\\
\lsg{(c)} Die Drehrichtung bleibt gleich.\\
\textit{Begründung: Beide Änderungen heben sich auf. Strom umgekehrt + B-Feld umgekehrt = Kraft bleibt gleich.}

\end{successbox}

% ═══════════════════════════════════════════════════════════════
% 9. LÖSUNGEN LB-09a
% ═══════════════════════════════════════════════════════════════

\section*{9. Lösungen AB-\stession}

\begin{successbox}[Lückentext-Lösungen]

\textbf{Kasten 1: Aufbau}
\begin{itemize}
    \item Stator = \lsg{feststehender Teil}
    \item Rotor = \lsg{drehbarer Teil}
    \item Kommutator = \lsg{Polwender / Stromwender}
    \item Bürsten = \lsg{Schleifkontakte}
\end{itemize}

\textbf{Kasten 2: Funktionsprinzip}
\begin{enumerate}
    \item Strom fließt durch die \lsg{Spule}
    \item Im \lsg{Magnetfeld} des Stators
    \item Die \lsg{Lorentzkraft} wirkt
    \item Die Spule beginnt sich zu \lsg{drehen}
    \item Nach \lsg{180}° würde Kraft zurückdrücken
    \item Der \lsg{Kommutator} kehrt Stromrichtung um
    \item Energieumwandlung: \lsg{Elektrische} $\rightarrow$ \lsg{Kinetische/mechanische}
\end{enumerate}

\textbf{Kasten 3: Drehrichtung ändern}
\begin{itemize}
    \item \lsg{Umpolen} der Spannung
    \item \lsg{Umpolen/Umdrehen} des Magnetfelds
    \item Beides gleichzeitig: Drehrichtung bleibt \lsg{gleich}
\end{itemize}

\textbf{Aufgabe 1 ($\star$, 4 BE):}
\begin{tabular}{lcc}
Aussage & R & F \\
Rotor ist feststehend & & \lsg{X} (dreht sich) \\
Kommutator kehrt Strom um & \lsg{X} & \\
Bürsten drehen sich mit & & \lsg{X} (feststehend) \\
Ohne Kommutator keine Drehung & & \lsg{X} (pendelt nur) \\
\end{tabular}

\textbf{Aufgabe 2 ($\star$, 3 BE):} \lsg{Staubsauger, Ventilator, Akkuschrauber} (Kreuze)

\textbf{Aufgabe 3 ($\star\star$, 2 BE):} \lsg{Ohne Kommutator würde der Rotor stehenbleiben (nach kurzem Pendeln), weil der N-Pol des Rotors vom S-Pol des Stators angezogen wird.}

\textbf{Aufgabe 4 ($\star\star\star$, 3 BE):} \lsg{Nein}, er hat nicht recht. Wenn man beides gleichzeitig ändert, heben sich die Änderungen auf. Die Kraft wirkt weiterhin in dieselbe Richtung.

\end{successbox}

% ═══════════════════════════════════════════════════════════════
% 10. AUSBLICK
% ═══════════════════════════════════════════════════════════════

\section*{10. Ausblick: Nächste Stunde (Induktion)}

\begin{infobox}[Überleitung zur Induktion]
\textbf{Leitfrage:} ``Der Motor wandelt elektrische Energie in Bewegung um. Geht das auch umgekehrt?''

\textbf{Demonstration (Einstieg nächste Stunde):}
\begin{itemize}
    \item Motor an Voltmeter anschließen (statt Stromquelle)
    \item Rotor von Hand drehen
    \item Spannung wird angezeigt $\rightarrow$ Generator!
\end{itemize}

\textbf{Merksatz:} Motor und Generator sind ``Geschwister'' -- gleicher Aufbau, umgekehrte Funktion.
\end{infobox}

\end{document}
